\paragraph{三成分閉包と局所モデル.}
一般の $Q$ について、$X_Q$ は滑らかな楕円線織面 $R$ と二つの平面
$\Pi_\pm\simeq\mathbb P^2$ からなる単純正規交差曲面である。
二つの共通等方三平面を $U_\pm$ とすると
$\Pi_\pm=\{U_\pm\oplus\ell:\ell\in\mathbb P(K_3)\}$ である。
射影 $W\to K_3$ の階数二は、Pfaffian三次曲線に直線を強制するため存在しない。
階数一の点はちょうどこの二平面、階数三の点は滑らかな一つの既約開部分となる。

$K_6=U_+\oplus U_-$ に対する交差ブロックを $M(z)$ とすれば、
境界近傍で全方程式は正確に
\[
 M(\ell)d=0
\]
へ消去できる。
$\det M=0$ は滑らかな三次曲線なので、局所式は $fv=0$ となる。
これは成分や交差を点集合だけでなくスキームとして決める小さなモデルである。
$R$ は $E$ 上の分解可能な射影直線束で、二重曲線 $D_\pm$ は互いに交わらない二つの切断となる。
境界平面の $\ell$ と線織面の底の $z$ を同じ座標と仮定しない。

\[
 \deg N_{D_\pm/R}=0,\qquad
 \deg N_{D_\pm/\Pi_\pm}=9,\qquad
 \deg\mathcal T_X^1|_{D_\pm}=9.
\]
一般一次摂動は種数二十K3曲面を与え、その全空間には二つの二重曲線上に九個ずつ普通二重点が生じる。
この九点は一次平滑化切断の単純零点である。
二つの平面成分のイデアルを吹き上げる局所式は $ab=tc$ の $(a,t)$ のblow-upであり、
射影的な小さなcrepant解消を与える。
中心ファイバーは、中央の $R$ を保ち、両側の平面をそれぞれ九点で吹き上げた三成分のType II Kulikovモデルとなる。
法束は平滑化零因子を引くことで相殺する。
九点を吹き上げた各平面を、追加条件なしに有理楕円曲面とは呼ばない。

\paragraph{スクロールの分極と再利用可能な数値.}
$\det\mathcal K=\mathcal O_E(-2)$ なので、
直線部分空間をパラメータ化する $\mathbb P_E(\mathcal K)$ の超平面類の平方は六である。
GrassmannのPl\"ucker類については $\deg X_Q=38$、両平面の次数は一ずつ、
主成分 $R$ の次数は三十六である。
これら二つの分極を同一視しない。
一般関数体上で二平面を区別するには、前述の二次拡大 $L/K$ が必要であり、
Kulikov修正の粗モジュライ上への降下とBrauer classの境界での分岐は残る問題である。
